% Einheit 2 von 3 – Die Nullhypothese wählen: aus der Sicht des Entscheiders (bank/hypothesentests/e2.jsonl, zusammenbau v0.1)
%% TODO Zeitmarke Einheit 2: Marken-Zeile nicht lesbar
%% TODO Zweigzeile Teil 1: was hier gelernt wird (Satz in Schülersprache aus dem Einheitstitel) – Einheit 2
\einheitenkopf[e2]{Einheit 2 von 3 · Die Nullhypothese wählen: aus der Sicht des Entscheiders}
\zweigzeile{Abitur LK}
\setcounter{aufgabe}{10}

%% TODO Titel: Ich-kann-Satz fehlt in der Bank (Platzhalter: Kettenname) – Nr. 11
%% TODO Anweisung: ein Satz über der ersten Teilaufgabe fehlt in der Bank – Nr. 11
\begin{aufgabe}{Wahl der Nullhypothese}
%% TODO \rechenplatz{n}: Schrittzahl des Grundfalls fehlt in der Bank (nur bei fester Schreibform, 2.3 b)
\begin{teile}
\teil $H_0$: „Das neue Medikament wirkt höchstens so oft wie das alte.“ In Wahrheit wirkt es nicht öfter, der Test lehnt $H_0$ aber ab. Welcher Fehler liegt vor? \\ \kreuz{Fehler erster Art} \\ \kreuz{Fehler zweiter Art} \\ \kreuz{kein Fehler}
\teil Bisher sind 5\,\% der Fahrradschläuche eines Herstellers undicht. Ein neues, teureres Gummi soll den Anteil senken. Getestet wird $H_0$: „Der Anteil undichter Schläuche beträgt mindestens 5\,\%.“ Gib eine Überlegung an, die zur Wahl dieser Nullhypothese geführt haben könnte, und begründe sie. \hfill (Abitur 2018 LK)
\teil Ein Verlag setzt bereits einen Algorithmus ein, der Artikel empfiehlt; bisher sind 65\,\% der Abonnenten zufrieden. Er soll nur abgeschaltet werden, wenn die Zufriedenheit nachweislich gesunken ist. Getestet wird $H_0$: „Mindestens 65\,\% der Abonnenten sind zufrieden.“ Gib eine mögliche Überlegung des Verlags zur Wahl dieser Nullhypothese an. \hfill (Abitur 2024 LK)
\teil Ein Händler hat mit einem Hersteller vereinbart, dass höchstens 5\,\% der gelieferten Ladegeräte defekt sein dürfen; sonst geht die Lieferung zurück. Er testet $H_0$: „Der Anteil defekter Geräte beträgt mindestens 5\,\%“ und nicht $H_0$: „Der Anteil beträgt höchstens 5\,\%.“ Gib eine Überlegung an, die zu dieser Wahl geführt haben könnte, und begründe sie. \hfill (Abitur 2022 LK)
\end{teile}
\end{aufgabe}

%% TODO Titel: Ich-kann-Satz fehlt in der Bank (Platzhalter: Kettenname) – Nr. 12
%% TODO Anweisung: ein Satz über der ersten Teilaufgabe fehlt in der Bank – Nr. 12
\begin{aufgabe}{Wahl der Nullhypothese – Fehler finden, Begründen, Anwendung}
\begin{teile}
\teil Mia begründet die Wahl von $H_0$: „Der Anteil defekter Geräte beträgt mindestens 5\,\%“ so: „Dann ist die Wahrscheinlichkeit, eine schlechte Lieferung anzunehmen, als Fehler zweiter Art begrenzt.“ Finde den Fehler.
\teil Begründe, warum nur der Fehler erster Art durch das Signifikanzniveau begrenzt ist.
\teil Eine Gemeinde will eine teure Lärmschutzwand nur bauen, wenn nachgewiesen ist, dass mehr als 25\,\% der Anwohner sich durch Verkehrslärm gestört fühlen. Wähle die Nullhypothese und begründe.
\end{teile}
\end{aufgabe}
